% === Recommendations ===
\section{How to Evaluate Incident Triage}
\label{sec:recommendations}

The findings above are only useful if they change practice. This section states
what we would ask of a future incident-triage evaluation, including our own.
Each recommendation is tied to a specific way our own results would have
misled, had we not run the corresponding check.

\smallskip
\noindent\textbf{R1. Report retrieval on a real corpus; do not let
classification stand alone.} Ticket-worthiness classification is near-perfectly
separable on injected faults (ROC area $1.00$) because an active fault differs
from a healthy baseline by construction. On real engineer-written tickets the
same question is only weakly answerable from text (ROC area $0.57$--$0.70$,
\S\ref{sec:results-rq4}). A strong classification number on fault-injection data
therefore certifies the dataset, not the method. Retrieval, by contrast, is
measured identically on both and transfers.

\smallskip
\noindent\textbf{R2. Compare trained components seed-matched, with seed-aware
intervals.} Retriever fine-tuning is the only stochastic step in our pipeline,
and it contributes roughly $28\%$ of \hit{5} variance on one dataset and $1\%$
on another (\S\ref{sec:results-seeds}). Our own headline margin was originally
computed from one run of each system against a query-bootstrap interval that
ignored seed variance entirely; the bi-encoder draw it happened to use sat near
the low end of that retriever's range, and the interval was roughly $2.6\times$
too narrow. Report $\sqrt{\sigma^2_{\text{seed}}+\sigma^2_{\text{sample}}}$, or
the paired per-seed difference, whenever a retrained component is involved.

Two corollaries. First, \emph{a fixed seed is not reproducibility} for GPU
training: re-running our own configuration unchanged moved \hit{5} by $0.007$.
Second, where the ranges of two systems do not overlap across seeds, say so ---
it is a stronger and more assumption-free claim than any difference of means.

\smallskip
\noindent\textbf{R3. When ground truth is inferred, report more than one
relation, and calibrate against chance.} Moving from our coarse relation to a
stricter one changes \hit{5} from $0.970$ to $0.744$ and reorders the method
table. Neither number is meaningful without knowing how much of memory each
relation admits: the average held-out query has $410$ relevant items under the
coarse relation and $50$ under the strict one, so a uniformly random ranker
would score $0.051$ and $0.006$ respectively. Publishing that reference costs
one sentence and lets a reader judge the result independently.

\smallskip
\noindent\textbf{R4. Audit the feature matrix before interpreting a
classifier.} Our own reported finding that real-incident triage sits at chance
came from classifiers fitted to $94$ telemetry feature columns that are
zero-filled for a text-only corpus --- $0$ of $78{,}140$ rows carried a
non-zero value. The result was a constant predictor, and its bootstrap interval
was exactly $[0.5000, 0.5000]$, which should have been read as a symptom rather
than a tight estimate. A one-line check for degenerate inputs would have caught
it; we now run one as the first output of that experiment.

\smallskip
\noindent\textbf{R5. Separate measured quantities from modelled ones in cost
claims.} Our cost saving is a counterfactual: measured invocation and
gate-skip counts, multiplied by modelled per-call costs. Calling it a measured
wall-clock reduction, as we originally did, overstates it. The honest version is
also more informative --- because every gated step sits behind the same gate,
the saving approximates that gate's closure rate almost independently of the
cost table, which we verify by recomputing with measured durations
($79.0\%$ versus $77.9\%$).

Note also what a naive fix would have done. Our evaluation harness replays
precomputed predictions, so its measured end-to-end time is $43$\,ms per
window: the cost of replaying a decision, not of making one. Reporting that as
a deployment figure would have been considerably worse than the estimate.

\smallskip
\noindent\textbf{R6. Evaluate at more than one corpus scale.} Our conclusion
about retrieval fusion \emph{reverses} between a $347$-item memory and a
$38{,}642$-item one: fusion is significantly worse than a single dense retriever
on the former and significantly better on the latter
(\S\ref{sec:results-rq1}). Had we evaluated only on the demonstration systems,
as is common, we would have reported that fusing retrievers is harmful. Scale
and heterogeneity are not incidental properties of a benchmark; they decide the
answer.

\smallskip
\noindent\textbf{R7. State which metric the claim is about.} Fusion improves
\hit{5} and simultaneously \emph{degrades} \hit{1}
(\S\ref{sec:results-rq2}). Both are significant; neither is the whole picture.
Which one matters depends on whether the consumer is an engineer scanning a
shortlist or an automated handoff acting on the top result. Reporting a single
headline metric hides the trade-off.
